% Generated by scripts/assemble.py; corrections belong in source/corrections.json.
\examhead{88}{苏州大学 2024 年数学分析试卷}{87}
% p088-b01; source PDF page 88
\par\Needspace{5\baselineskip}
\blockanchor{p088-b01}
\majorquestion{一}计算题. 每题 10 分, 共 20 分.\par
% p088-b02; source PDF page 88
\blockanchor{p088-b02}
\question{1}求极限 $\lim_{x\to 0}\frac{\sin x - \arctan x}{x^3}.$\par
% p088-b03; source PDF page 88
\blockanchor{p088-b03}
\question{2}设  $f(x)=x^{3}\cos x$ ，求  $f^{(2023)}(0)$ .\par
% p088-b04; source PDF page 88
\par\Needspace{5\baselineskip}
\blockanchor{p088-b04}
\majorquestion{二}(15 分) 设 $a_1 = 1$, 且当 $n \geq 2$ 时, $a_n = \frac{1}{2} \left( a_{n-1} + \frac{5}{a_{n-1}} \right)$. 证明: $\{a_n\}$ 收敛, 并求极限.\par
% p088-b05; source PDF page 88
\par\Needspace{5\baselineskip}
\blockanchor{p088-b05}
\majorquestion{三}(15 分) 设 $f(x) = \begin{cases} \frac{\sin x}{x}, & x \neq 0; \\ 1, & x = 0. \end{cases}$ 证明: 函数 $f(x)$ 在 $(-\infty, +\infty)$ 上一致连续.\par
% p088-b06; source PDF page 88
\par\Needspace{5\baselineskip}
\blockanchor{p088-b06}
\majorquestion{四}(15 分) 计算三重积分 $\iiint_{\Omega} (x + y + z)^2 \, \mathrm{d}x \, \mathrm{d}y \, \mathrm{d}z$, 其中 $\Omega = \{(x, y, z) \in \mathbb{R}^3 \mid |x| + |y| + |z| \leq 1\}$.\par
% p088-b07; source PDF page 88
\par\Needspace{5\baselineskip}
\par\Needspace{6\baselineskip}
\blockanchor{p088-b07}
\majorquestion{五}（15分）设 $u = f(x,y,z)$ 的所有二阶偏导数都连续，$v = f(x(s,t,r),y(s,t,r),z(s,t,r))$，其中\par
\nopagebreak[4]
% p088-b08; source PDF page 88
\blockanchor{p088-b08}
\displayformula{\left\{ \begin{array}{l} x (s, t, r) = \frac {1}{9} (a x + 4 t + 8 r); \\ y (s, t, r) = \frac {1}{9} (4 s + b t - 4 r); \\ z (s, t, r) = \frac {1}{9} (8 s - 4 t + c r). \end{array} \right.}
% p088-b09; source PDF page 88
\par\Needspace{6\baselineskip}
\blockanchor{p088-b09}
试讨论是否存在常数 $a, b, c$，使得当 $x = x(s, t, r), y = y(s, t, r), z = z(s, t, r)$ 时，总成立\par
\nopagebreak[4]
% p088-b10; source PDF page 88
\blockanchor{p088-b10}
\displayformula{\left. \left(u _ {x x} + u _ {y y} + u _ {z z}\right) \right| _ {(x, y, z)} = \left. \left(v _ {s s} + v _ {t t} + v _ {r r}\right) \right| _ {s, t, r}.}
% p088-b11; source PDF page 88
\blockanchor{p088-b11}
若存在, 求 a, b, c 的值.\par
% p088-b12; source PDF page 88
\par\Needspace{5\baselineskip}
\blockanchor{p088-b12}
\majorquestion{六}（10分）设 $u(x,y,z) = \frac{x^2}{2} +\frac{y^2}{3} +\frac{z^2}{4}$ 求在原点处函数 $u$ 增长最快的方向.\par
% p088-b13; source PDF page 88
\par\Needspace{5\baselineskip}
\blockanchor{p088-b13}
\majorquestion{七}(10 分) 证明不等式: $\int_{0}^{\frac{\pi}{2}} \frac{\sin x}{1 + x^{2}} \, \mathrm{d}x < \int_{0}^{\frac{\pi}{2}} \frac{\cos x}{1 + x^{2}} \, \mathrm{d}x$.\par
% p088-b14; source PDF page 88
\par\Needspace{5\baselineskip}
\blockanchor{p088-b14}
\majorquestion{八}(10 分) 设函数 $f$ 在 $[0,1]$ 上二阶可导, $f(0) = 0$, $f(1) = 1$, 并存在 $c \in (0,1)$, 使得 $f(c) > c^2$. 证明: 至少存在 $\xi \in (0,1)$, 使得 $f''(\xi) < 2$.\par
% p088-b15; source PDF page 88
\par\Needspace{5\baselineskip}
\blockanchor{p088-b15}
\majorquestion{九}(10 分) 设 $f$ 是 $[0, +\infty)$ 上无界的连续函数. 问: $\int_0^{+\infty} f(x) \mathrm{d}x$ 是否发散? 给出证明或反例.\par
% p088-b16; source PDF page 88
\par\Needspace{5\baselineskip}
\blockanchor{p088-b16}
\majorquestion{十}(10 分) 设 $\{f_n(x)\}$ 是 $[a, b]$ 上黎曼可积的函数列, 且 $\{f_n(x)\}$ 在 $[a, b]$ 上一致收敛于 $f$. 证明:\par
% p088-b17; source PDF page 88
\blockanchor{p088-b17}
(1) $f$ 在 $[a, b]$ 上可积.\par
% p088-b18; source PDF page 88
\blockanchor{p088-b18}
(2) $\{F_n(x)\}$ 在 $[a, b]$ 上一致收敛于 $F(x)$, 其中 $F_n(x) = \int_{a}^{x} f_n(t) \mathrm{d}t$, $F(x) = \int_{0}^{x} f(t) \mathrm{d}t$.\par
% p088-b19; source PDF page 88
\par\Needspace{5\baselineskip}
\blockanchor{p088-b19}
\majorquestion{十一}(10 分) 设 $u$ 是平面开区域 $D$ 上的二元函数, 且所有的偏导数连续. 证明: $u$ 是 $D$ 上的调和函数, 即在 $D$ 上 $\frac{\partial^2 u}{\partial x^2} + \frac{\partial^2 u}{\partial y^2} \equiv 0 \Longleftrightarrow$ 对 $D$ 内任意圆周 $L$, 有 $\oint_L \frac{\partial u}{\partial n} = 0$, 其中 $\frac{\partial u}{\partial n}$ 表示 $u$ 在 $L$ 上的外法向导数.\par
% p088-b20; source PDF page 88
\par\Needspace{5\baselineskip}
\blockanchor{p088-b20}
\majorquestion{十二}(10 分) 证明: 存在 $\xi \in (0, \pi)$, 使得 $\sum_{n=1}^{\infty} \frac{\cos n\xi}{n \ln^{2}(n+2)} = 0$.\par
